% ID: FIS-CIR-OHM-001
% Titolo: Relazione tra tensione, corrente e resistenza
% Disciplina: Fisica
% Area: Circuiti elettrici
% Argomento: Legge di Ohm
% Sottoargomento: Relazione tra grandezze elettriche
% Classe: Prima istituto professionale
% Difficolta: 2
% Tipo: quesito_teorico
% Risultato: V = R I
% Tempo_stimato: 5 min
% Competenze: formalizzazione matematica, collegamento tra grandezze fisiche, uso di una legge
% Prerequisiti: tensione, corrente elettrica, resistenza
% Tag: fisica, circuiti_elettrici, legge_ohm, tensione, corrente, resistenza
% Fonte: verifica_fisica_circuiti_anonimizzata
% Autore: Riccardo Carli
% Licenza: CC-BY-SA-4.0

\begin{esercizio}
Come sono legate tra loro le grandezze fisiche resistenza, intensita di
corrente e tensione? Scrivi la relazione matematica corrispondente.
\end{esercizio}

\begin{soluzione}
Per un conduttore ohmico tensione, corrente e resistenza sono legate dalla
legge di Ohm:
\[
V=RI.
\]

La tensione \(V\) si misura in volt, la corrente \(I\) in ampere e la
resistenza \(R\) in ohm. Dalla stessa relazione si possono ricavare anche:
\[
I=\frac{V}{R}
\qquad\text{e}\qquad
R=\frac{V}{I}.
\]

Quindi, a resistenza fissata, se la tensione aumenta aumenta anche la corrente;
a tensione fissata, se la resistenza aumenta la corrente diminuisce.
\end{soluzione}
